% ══════════════════════════════════════════════════════════════
%  Section 9.4: UTXO Set --- The Spendability Database
% ══════════════════════════════════════════════════════════════
\chapter{Chain: UTXO Set}
\label{ch:utxo-set}

\begin{learnbox}
\begin{itemize}
  \item What a UTXO (Unspent Transaction Output) is and why Bitcoin tracks outputs, not accounts
  \item How the UTXO set is persisted in sled: one tree with \code{txid\_bytes -> Vec<TXOutput>}
  \item Coin selection: how \code{find\_spendable\_outputs} chooses which coins to spend
  \item Block connection: updating the UTXO set when a block is accepted (remove spent, add created)
  \item Chain reorg support: rolling back the UTXO set when a block is disconnected
  \item Balance queries: deriving ``account balance'' from the collection of spendable UTXOs
\end{itemize}
\end{learnbox}

In Bitcoin, a \textbf{\index{UTXO}UTXO} (Unspent \index{transaction}Transaction Output) is a \emph{specific coin} you can spend: an output created by some earlier \index{transaction}transaction that has not yet been referenced as an input by any later \index{transaction}transaction. Your ``balance'' is therefore not a single stored number; it is the sum of the values of all \index{UTXO}UTXOs you can unlock with your keys.

The \textbf{\index{UTXO}UTXO set} is the node's \emph{current index of spendability}: the collection of all \index{UTXO}UTXOs that exist \textbf{right now} (keyed by \(txid, vout\) and mapping to output details). It is updated when \index{block}blocks are connected (remove spent outputs, add newly created outputs), and it is what \index{wallet}wallets and validators consult to answer ``is this input spendable?'' efficiently.

This is different from the \textbf{\index{blockchain}blockchain}, which is an append-only history. The chain records \emph{what happened}; the \index{UTXO}UTXO set represents the \emph{resulting state} (what remains spendable after applying that history).

% ──────────────────────────────────────────────────────────────
\section{Storage Layout in sled}
\label{sec:ch09-4-storage-layout}

\index{sled}\code{sled} is an embedded key-value database written in Rust. There is no \code{\index{UTXO}UTXO} struct. The \index{UTXO}UTXO set is stored as a \index{sled}sled tree:

\begin{itemize}
  \item Tree name: \code{chainstate} (constant \code{\index{UTXO}UTXO\_TREE})
  \item Key: \code{\index{transaction}txid} (raw bytes---the 32-byte \index{transaction}transaction hash)
  \item Value: \code{Vec<TXOutput>} serialized via \code{Rust!serialization}\code{bincode}
\end{itemize}

\begin{listing}[htbp]
\caption{UTXO tree structure and handle}
\label{lst:ch09-4-1}
\begin{minted}[fontsize=\footnotesize]{rust}
// Source: utxo_set.rs
const UTXO_TREE: &str = "chainstate";

pub struct UTXOSet {
    // handle used to access the DB and chain data
    blockchain: BlockchainService,
}
\end{minted}
\end{listing}

\begin{notebox}
We persist \textbf{derived state} as: \code{txid\_bytes -> bincode(Vec<TXOutput>)}. The \emph{outpoint} \((txid, vout)\) is interpreted as ``the output at index \code{vout} inside that stored \code{Vec<TXOutput>}''.

We use bytes as the database key because it is the \textbf{canonical representation}, more space/time efficient, and avoids repeated hex encoding/decoding. When a transaction ID is displayed as hex in logs, that is the same value, just rendered for humans.
\end{notebox}

% ──────────────────────────────────────────────────────────────
\section{Core Data Types: TXInput and TXOutput}
\label{sec:ch09-4-datatypes}

Before we can understand the UTXO set, we need to understand what a transaction \emph{points at} and what an output \emph{contains}.

\begin{listing}[htbp]
\caption{\code{TXInput}: reference to a previous output and authorization}
\label{lst:ch09-4-2}
\begin{minted}[fontsize=\footnotesize]{rust}
// Source: transaction.rs
#[derive(Clone, Default, Serialize, Deserialize)]
pub struct TXInput {
    txid: Vec<u8>,                  // previous transaction hash
    vout: usize,                    // output index in that transaction
    signature: Vec<u8>,             // proof of authorization
    pub_key: Vec<u8>,               // public key for verification
}

impl TXInput {
    pub fn new(txid: &[u8], vout: usize) -> TXInput {
        TXInput {
            txid: txid.to_vec(),
            vout,
            signature: vec![],
            pub_key: vec![],
        }
    }

    pub fn uses_key(&self, pub_key_hash: &[u8]) -> bool {
        let locking_hash = hash_pub_key(self.pub_key.as_slice());
        locking_hash.eq(pub_key_hash)
    }
}
\end{minted}
\end{listing}

\begin{notebox}
An \textbf{outpoint} is the pair \((txid, vout)\) that identifies which previous output an input is spending. It does not store ``which address''; it stores the exact transaction and output index being consumed.
\end{notebox}

\begin{listing}[htbp]
\caption{\code{TXOutput}: value and the ``in mempool'' flag}
\label{lst:ch09-4-3}
\begin{minted}[fontsize=\footnotesize]{rust}
// Source: transaction.rs
#[derive(Clone, Serialize, Deserialize)]
pub struct TXOutput {
    value: i32,                       // amount locked by this output
    in_global_mem_pool: bool,         // local hint: reserved in mempool?
    pub_key_hash: Vec<u8>,            // ``lock'': who can spend this
}

impl TXOutput {
    pub fn new(
        value: i32,
        address: &WalletAddress
    ) -> Result<TXOutput> {
        let mut output = TXOutput {
            value,
            in_global_mem_pool: false,
            pub_key_hash: vec![],
        };
        output.lock(address)?;
        Ok(output)
    }

    fn lock(&mut self, address: &WalletAddress) -> Result<()> {
        self.pub_key_hash = get_pub_key_hash(address)?;
        Ok(())
    }

    pub fn set_in_global_mem_pool(&mut self, value: bool) {
        self.in_global_mem_pool = value;
    }
}
\end{minted}
\end{listing}

\begin{notebox}
\code{pub\_key\_hash} is the ``lock''---it is what \code{is\_locked\_with\_key()} matches during coin selection. \code{in\_global\_mem\_pool} is \textbf{not} a Bitcoin consensus concept; it is a \emph{local hint} in this learning implementation used to avoid selecting an output already referenced in a pending transaction.
\end{notebox}

% ──────────────────────────────────────────────────────────────
\section{Coin Selection: Finding Spendable Outputs}
\label{sec:ch09-4-coin-selection}

When creating a transaction, the wallet needs to answer: ``Which outpoints locked to me can cover \(\mathit{amount}\), while avoiding outputs already used in the mempool?''

\begin{listing}[htbp]
\caption{\code{UTXOSet::find\_spendable\_outputs} --- scan and select}
\label{lst:ch09-4-4}
\begin{minted}[fontsize=\footnotesize]{rust}
// Source: utxo_set.rs
pub async fn find_spendable_outputs(
    &self,
    from_pub_key_hash: &[u8],
    amount: i32,
) -> Result<(i32, HashMap<String, Vec<usize>>)> {
    let mut unspent_outputs_indexes = HashMap::new();
    let mut accumulated = 0;
    let db = self.blockchain.get_db().await?;
    let utxo_tree = db.open_tree(UTXO_TREE)?;

    for item in utxo_tree.iter() {
        let (k, v) = item?;
        let txid_hex = HEXLOWER.encode(k.as_ref());
        let (tx_out, _) = bincode::serde::decode_from_slice(
            v.as_ref(),
            bincode::config::standard(),
        )?;

        for (vout_idx, out) in tx_out.iter().enumerate() {
            if out.not_in_global_mem_pool()
                && out.is_locked_with_key(from_pub_key_hash)
                && accumulated < amount
            {
                accumulated += out.get_value();
                unspent_outputs_indexes
                    .entry(txid_hex.clone())
                    .or_insert_with(Vec::new)
                    .push(vout_idx);
            }
        }
    }
    Ok((accumulated, unspent_outputs_indexes))
}
\end{minted}
\end{listing}

\begin{checkpointbox}
This method:
\begin{itemize}
  \item Opens the \code{UTXO\_TREE} and iterates the entire keyspace (this is \(O(N)\) over all stored outputs)
  \item Deserializes each \code{Vec<TXOutput>} and checks each output for ownership and reservation status
  \item Accumulates outputs until the sum covers the requested amount
  \item Returns a \code{HashMap<txid\_hex -> [vout indexes]>} mapping transaction IDs to their selected output indices
\end{itemize}
\end{checkpointbox}

% ──────────────────────────────────────────────────────────────
\section{Transaction Construction Using Selected Outputs}
\label{sec:ch09-4-tx-construction}

Coin selection results are used by \code{Transaction::new\_utxo\_transaction} to build the actual transaction:

\begin{listing}[htbp]
\caption{UTXO transaction construction}
\label{lst:ch09-4-5}
\begin{minted}[fontsize=\footnotesize]{rust}
// Source: transaction.rs
pub async fn new_utxo_transaction(
    from_wlt_addr: &WalletAddress,
    to_wlt_addr: &WalletAddress,
    tx_amount: i32,
    utxo_set: &UTXOSet,
) -> Result<Transaction> {
    let wallets = WalletService::new()?;
    let from_wallet = wallets.get_wallet(from_wlt_addr)?;
    let from_pub_key_hash = hash_pub_key(
        from_wallet.get_public_key()
    );
    let (available_funds, valid_outputs) = utxo_set
        .find_spendable_outputs(&from_pub_key_hash, tx_amount)
        .await?;

    if available_funds < tx_amount {
        return Err(BtcError::NotEnoughFunds);
    }

    let mut inputs = vec![];
    for (txid_hex, out_indexes) in valid_outputs {
        let txid = HEXLOWER.decode(txid_hex.as_bytes())?;
        for vout_idx in out_indexes {
            inputs.push(TXInput::new(&txid, vout_idx));
        }
    }

    let mut outputs = vec![TXOutput::new(tx_amount, to_wlt_addr)?];
    if available_funds > tx_amount {
        outputs.push(TXOutput::new(
            available_funds - tx_amount,
            from_wlt_addr,
        )?);
    }

    let mut tx = Transaction {
        id: vec![],
        vin: inputs,
        vout: outputs,
    };
    tx.id = tx.hash()?;
    tx.sign(utxo_set.get_blockchain(), from_wallet.get_pkcs8())
        .await?;
    Ok(tx)
}
\end{minted}
\end{listing}

\begin{notebox}
The selected outpoints become \code{TXInput}s, value is re-expressed as payment and optional change \code{TXOutput}s, then we compute a txid and sign. This is the \textbf{mechanical implementation of the UTXO model}: spending consumes entire outputs, so leftover value must be recreated as a new change output.
\end{notebox}

% ──────────────────────────────────────────────────────────────
\section{Updating After a Block (Block Connect)}
\label{sec:ch09-4-block-connect}

When a block is accepted, the UTXO set must be updated: spent outputs are removed, newly created outputs are added.

\begin{listing}[htbp]
\caption{Update the UTXO set when a block is connected}
\label{lst:ch09-4-6}
\begin{minted}[fontsize=\footnotesize]{rust}
// Source: utxo_set.rs
pub async fn update(&self, block: &Block) -> Result<()> {
    let db = self.blockchain.get_db().await?;
    let utxo_tree = db.open_tree(UTXO_TREE)?;

    // Step 1: Remove (spend) inputs referenced by this block
    for tx in block.get_transactions() {
        if !tx.is_coinbase() {
            for input in tx.get_vin() {
                let txid = input.get_txid();
                if let Ok(Some(mut outputs)) =
                    self.find_utxo(txid)
                {
                    if input.vout < outputs.len() {
                        outputs.remove(input.vout);
                        if outputs.is_empty() {
                            utxo_tree.remove(txid)?;
                        } else {
                            let data = bincode::serde::encode_to_vec(
                                &outputs,
                                bincode::config::standard(),
                            )?;
                            utxo_tree.insert(txid, data)?;
                        }
                    }
                }
            }
        }
    }

    // Step 2: Add (create) outputs produced by this block
    for tx in block.get_transactions() {
        let txid = tx.get_id();
        let outputs = tx.get_vout();
        let data = bincode::serde::encode_to_vec(
            &outputs,
            bincode::config::standard(),
        )?;
        utxo_tree.insert(txid, data)?;
    }

    Ok(())
}
\end{minted}
\end{listing}

\begin{notebox}
Block connection is a two-step atomic transformation:
\begin{enumerate}
  \item \textbf{Spend phase}: for each input in each transaction, remove that output from the UTXO set (delete the txid key if all outputs are spent)
  \item \textbf{Create phase}: for each transaction in the block, insert all its outputs into the UTXO set (keyed by the transaction's ID)
\end{enumerate}
\end{notebox}

% ──────────────────────────────────────────────────────────────
\section{Rollback: Disconnecting a Block (Reorg Support)}
\label{sec:ch09-4-rollback}

If a block must be disconnected (e.g., during a chain reorg), we undo its UTXO effects: remove created outputs, restore spent outputs.

\begin{listing}[htbp]
\caption{Rollback the UTXO set when a block is disconnected}
\label{lst:ch09-4-7}
\begin{minted}[fontsize=\footnotesize]{rust}
// Source: utxo_set.rs
pub async fn rollback_block(&self, block: &Block) -> Result<()> {
    let db = self.blockchain.get_db().await?;
    let utxo_tree = db.open_tree(UTXO_TREE)?;

    // Step 1: Remove outputs created by this block
    for tx in block.get_transactions() {
        utxo_tree.remove(tx.get_id())?;
    }

    // Step 2: Restore outputs spent by this block
    for tx in block.get_transactions() {
        if !tx.is_coinbase() {
            for input in tx.get_vin() {
                let prev_txid = input.get_txid();
                if let Ok(Some(prev_tx)) =
                    self.blockchain.find_transaction(prev_txid)
                {
                    let data = bincode::serde::encode_to_vec(
                        prev_tx.get_vout(),
                        bincode::config::standard(),
                    )?;
                    utxo_tree.insert(prev_txid, data)?;
                }
            }
        }
    }

    Ok(())
}
\end{minted}
\end{listing}

\begin{notebox}
Rollback is the inverse of block connect:
\begin{enumerate}
  \item \textbf{Remove phase}: delete all transactions created by this block from the UTXO set
  \item \textbf{Restore phase}: for each input spent by this block, look up the original transaction and restore its full output set to the UTXO tree
\end{enumerate}

This allows the node to support chain reorganizations: if a better chain is discovered, blocks can be disconnected from the old chain and reconnected to the new one, with UTXO state updated correctly each step.
\end{notebox}

% ──────────────────────────────────────────────────────────────
\section{Balance and UTXO Queries}
\label{sec:ch09-4-queries}

Beyond the core update/rollback loop, we provide higher-level query methods for wallets and debugging.

\begin{listing}[htbp]
\caption{Find UTXOs and compute balance}
\label{lst:ch09-4-8}
\begin{minted}[fontsize=\footnotesize]{rust}
// Source: utxo_set.rs
pub async fn find_utxo(
    &self,
    pub_key_hash: &[u8],
) -> Result<Vec<TXOutput>> {
    let mut result = Vec::new();
    let db = self.blockchain.get_db().await?;
    let utxo_tree = db.open_tree(UTXO_TREE)?;

    for item in utxo_tree.iter() {
        let (_, v) = item?;
        let (tx_out, _) = bincode::serde::decode_from_slice(
            v.as_ref(),
            bincode::config::standard(),
        )?;
        for output in tx_out {
            if output.is_locked_with_key(pub_key_hash) {
                result.push(output);
            }
        }
    }
    Ok(result)
}

pub async fn get_balance(
    &self,
    address: &WalletAddress,
) -> Result<i32> {
    let pub_key_hash = get_pub_key_hash(address)?;
    let utxos = self.find_utxo(&pub_key_hash).await?;
    Ok(utxos.iter().map(|o| o.get_value()).sum())
}

pub async fn utxo_count(
    &self,
    address: &WalletAddress,
) -> Result<usize> {
    let pub_key_hash = get_pub_key_hash(address)?;
    let utxos = self.find_utxo(&pub_key_hash).await?;
    Ok(utxos.len())
}
\end{minted}
\end{listing}

\begin{notebox}
These query methods iterate the UTXO tree and filter by ownership (\code{pub\_key\_hash}). They answer:
\begin{itemize}
  \item ``What UTXOs do I own?'' (\code{find\_utxo})
  \item ``What is my total balance?'' (\code{get\_balance})
  \item ``How many unspent coins do I have?'' (\code{utxo\_count})
\end{itemize}
\end{notebox}

% ──────────────────────────────────────────────────────────────
\section{Summary}
\label{sec:ch09-4-summary}

\begin{chaptersummary}
\item \textbf{UTXO set} is the derived state that answers ``what is spendable right now?''
\item \textbf{Storage}: sled tree \code{chainstate} with \code{txid\_bytes -> Vec<TXOutput>}
\item \textbf{Coin selection}: \code{find\_spendable\_outputs} scans the UTXO tree, filtering by ownership and reservation status
\item \textbf{Block connect}: two-phase update (spend inputs, create outputs)
\item \textbf{Block disconnect}: two-phase rollback (remove created, restore spent) for reorg support
\item \textbf{Balance queries}: derive wallet balance by summing owned UTXOs
\item \textbf{Local mempool flag}: \code{in\_global\_mem\_pool} is a hint to avoid local double-selection, not a consensus rule
\end{chaptersummary}
